%%原模板参考了https://www.wondercv.com/的模板 说老实话 超级简历真比overleaf好用  就像低代码平台一样   但是overleaf使用latex语法支持更好的扩展性
\documentclass[11pt]{article}
% disable indent globally
\setlength{\parindent}{0pt}
% some general improvements, defines the XeTeX logo
\usepackage{xltxtra}
\usepackage{bookmark}
% use hyperlink for email and url
\usepackage{hyperref}
\hypersetup{hidelinks}
\usepackage{url}
\urlstyle{tt}
\usepackage{multicol}
\usepackage{xcolor}
%%%% 统一一种颜色，偏蓝色，用于section下划线和fontawesome
\definecolor{CVBlue}{RGB}{078,164,239 }
%   RGB值   可以关注微信公众号 阿昆的科研日常   那里有很多配色 我列一下我这个二次元喜欢的配色 (R G  B )   其他可以去看配色方案下面的文件
% 胡桃红    201 071  055  
% 芙芙蓝    078,164,239 
% 宵宫红    205 068  050
% 甘雨蓝    196 216  242
% 流萤绿    059 165  149 
% 电兵蓝   我才不找呢   哼哼哼




%照片   标准一寸白底照片   找照相馆拍   拍的好看一点    192*240 
%  姓名   作为简历开头大标题
%  给出学校标识   
% 给出充足的联系方式 
% 给出博客地址   base   很快入职 + 干半年

\usepackage{calc}
%%%% 利用tikz来定位照片和学校Logo 
\usepackage{graphicx}
\usepackage{tikz}
\usetikzlibrary{calc}
% loading fonts
\usepackage{fontspec}
\usepackage{xeCJK}
\CJKsetecglue{}
%% 取消中文与数字之间间隙
%%%%% 字体需要自己下载安装，注意版权问题。
% 这两种字体应该比较好看，英文Helvetica，中文方正兰亭黑，也是有多种版本，自己试试哪些好看。参考了https://www.wondercv.com/的模板
%%%%% windows系统好像需要先安装字体，之后下面语句就够了
%Main document font
%\setmainfont[
%  BoldFont = HelveticaNeueLTPro-Md.otf ,
%]{HelveticaNeueLTPro-Roman.otf}
%
%\setCJKmainfont[
%BoldFont=Pro_GB18030 DemiBold.otf,
%]{Pro_GB18030.otf}
%%%%% 字体需要自己下载安装，注意版权问题
%%%%% linux系统只需要字体路径就行了，如下
% % Main document font

\setmainfont[
Path = Font/,
  Extension = .otf ,
  BoldFont = HelveticaNeueLTPro-Md.otf ,
]{HelveticaNeueLTPro-Roman.otf}
\setCJKmainfont[
Path = Font/,
  Extension = .otf ,
BoldFont=ProGB18030 DemiBold.otf,
]{ProGB18030.otf}


\usepackage{relsize}
\usepackage{xspace}
\protected\def\Cpp{{C\nolinebreak[4]\hspace{-.05em}\raisebox{.28ex}{\relsize{-1}++}}\xspace} 
% use fontawesome
\usepackage{fontawesome}
%\newfontfamily{\FA}{[FontAwesome.otf]}
\usepackage[
	a4paper,
	left=1.2cm,
	right=1.2cm,
	top=1.5cm,
	bottom=1cm,
	nohead
]{geometry}
\renewcommand{\baselinestretch}{1.2} 
%定义行间距1.2
\usepackage{titlesec}
\usepackage{enumitem}
\setlist{noitemsep}

% removes spacing from items but leaves space around the whole list
%\setlist{nosep} % removes all vertical spacing within and around the list
\setlist[itemize]{topsep=0.25em, leftmargin=*}
\setlist[enumerate]{topsep=0.25em, leftmargin=*}
\titleformat{\section}        
% Customise the \section command 
  {\large\bfseries\raggedright} 
  % Make the \section headers large (\Large), % small capitals (\scshape) and left aligned (\raggedright)
  {}{0em}                      % Can be used to give a prefix to all sections, like 'Section ...'
  {}                           % Can be used to insert code before the heading
  [{\color{CVBlue}\titlerule}]                 % Inserts a horizontal line after the heading
\titlespacing*{\section}{0cm}{*1.6}{*1.2}
\usepackage{siunitx}
\usepackage{amssymb}
%\xeCJKsetup{CJKspace=true}
%\xeCJKDeclareCharClass{CJK}{`0 -> `9}    % 设置 0-9 以 CJK 字体输出
%\normalspacedchars{0,1,2,3,4,5,6,7,8,9} % 0-9 的字符类被还原


\begin{document}

\pagenumbering{gobble} 
%%%% 利用tikz来定位照片，标准一寸照片  192*240     985   
% \begin{tikzpicture}[remember picture, overlay]
% 	\node[anchor = north east] at ($(current page.north east)+(-1cm,-1.2cm)$) {\includegraphics[height=2.8cm]{img/myavator.jpg}};
% \end{tikzpicture}%
%%%% 利用tikz来定位学校Logo，这里只在第一页显示，345 × 77 mm 
\begin{tikzpicture}[remember picture, overlay]
	\node[anchor = north west] at ($(current page.north west)+(0.2cm,-0.2cm)$) {\includegraphics[height=2.5cm]{img/tongji.png}};
\end{tikzpicture}%
%%%% 利用tikz来定位页脚栏，电子版简历使用，黑白纸质打印效果可能并不好。这里只在第一页显示，如果需要每页都有，页脚或者background中加入。
% \begin{tikzpicture}[remember picture, overlay] 
%     \node[anchor = south,fill=CVBlue,draw=none,minimum width=\paperwidth,minimum height=1.5em,align=center,font=\footnotesize,text=white] at ($(current page.south)$) 
%     {\faGithubAlt \ \href{}{https://www.yuque.com/爱玩原神的可莉/}\qquad
%         \faRssSquare \ \href{github地址}{https://github.com/AHUA-Official} };
%         % 围栏填上自己的博客地址 
% \end{tikzpicture}

%tikzpicture环境很敏感，注释周围的空格、空行都会引起水平距离或垂直距离的变化，
%
\centerline{\LARGE\bfseries{罗力}}

% 联系信息+个人博客+求职优势(就是稳定干半年+入职快)
% 邮箱最好用你的edu邮箱   但是其实用qq也行  主要是要收到邮件查看响应的速度要快
\centerline{\normalsize{
		\faPhone \ 18228979484 \quad
		\faEnvelopeO \ \href{mailto:mxd@mail.bnu.edu.cn}{3167937401@qq.com}}}

\centerline{\normalsize{wx：ll18228979484 \quad github : https://github.com/No-99-Tongji }}
% 把你github地址写上去  

% 教育背景
\section{\makebox[\widthof{\faGraduationCap}][c]{\color{CVBlue}\faGraduationCap}\  教育背景}
\textbf{同济大学 \quad 软件工程 } \hfill 2023年9月-2027年6月 \\
平均绩点：4.54/5.0，专业排名前25\%，曾获校级二等奖学金，学习成绩优异。
% \newline {CET6}

% 我这个地方写太空了好像  

% % 实习经历   这个真的要好好写
\section{\makebox[\widthof{\faGraduationCap}][c]{\color{CVBlue}\faUsers}\ 实习经历}
\textbf{上海乐言科技股份有限公司 \hfill 后端开发实习生 \hfill 2026.1--至今}
\begin{itemize}
    \item \textbf{项目：} Hamilton系统 - 电商平台统一API代理与智能缓存系统
    \begin{itemize}
        \item \textbf{项目背景：} 统一对接抖音、淘宝、京东、拼多多等主流电商平台API，解决高频订单查询与API按次付费成本高的矛盾
        \item \textbf{核心目标：} 设计智能缓存系统，降低API调用成本，保障数据时效性
    \end{itemize}

    \item \textbf{技术方案：}
    \begin{itemize}
        \item \textbf{懒加载：} 上游请求触发更新，避免创建订单时的冗余API调用
        \item \textbf{缓存：} PostgreSQL持久化订单数据，为上游提供高效查询服务
        \item \textbf{定期更新：} Quarkus定时任务每秒扫描需更新店铺（查找last\_pulled\_at字段早于当前时间-阈值（如10分钟）的店铺），批量拉取这些店铺订单的最新状态
            \begin{itemize}
                \item \textbf{拉取间隙优化：} 10分钟间隙平衡时效性与成本：间隙太短导致冗余拉取（无状态更新），太长则时效性差；基于订单状态变化波峰分析（波峰集中在10分钟内和一天外），一天外波峰可被批量拉取覆盖，10分钟内波峰因拉取成本考虑不做高频更新
                \item \textbf{成本效益：} 通过优化拉取间隙，使API调用成本降至原来的20\%（成本公式：单次API调用成本 $\times$ QPS）
            \end{itemize}
        \item \textbf{定点更新：} 上游检测缓存过期时，触发单个订单实时更新% 检测机制待完善
    \end{itemize}

    \item \textbf{技术实现细节：}
    \begin{itemize}
        \item \textbf{数据库设计：} 在seller\_id和last\_pulled\_at字段建立索引，优化批量查询性能
        \item \textbf{系统架构：} 采用Controller-Service-Fetcher-Model分层架构
        \item \textbf{平台适配：} 参与拼多多、淘宝、抖音电商等异构API拉取器的搭建与维护
    \end{itemize}

    \item \textbf{技术选型与架构决策：}
    \begin{itemize}
        \item \textbf{Kotlin：} 公司原有Java技术栈基础上，选择Kotlin以获得更高性能，同时保持与Java模块的良好兼容性
        \item \textbf{Quarkus：} 采用Quarkus框架支持GraalVM原生编译，显著提升启动速度与运行效率，适合高频API代理场景
        \item \textbf{OpenAPI Generator：} 自动生成与各电商平台API对接的客户端代码，减少手动编写工作量，降低人为错误
        \item \textbf{架构设计：} Controller-Service-Fetcher-Model分层架构，清晰分离关注点，便于维护与扩展
    \end{itemize}

    \item \textbf{技术栈：} Kotlin, Quarkus, PostgreSQL, OpenAPI Generator
\end{itemize}


% 项目经历   这个真的要好好写
\section{\makebox[\widthof{\faGraduationCap}][c]{\color{CVBlue}\faUsers}\ 项目经历}
\textbf{EduNexus 学习管理系统 \hfill 组长/全栈开发 \hfill 2025.10--2025.12}
\begin{itemize}
    \item \textbf{团队与角色：} 3人全栈开发团队组长，负责技术选型、核心模块开发、团队协调与Scrum敏捷开发管理

    \item \textbf{技术选型与架构决策：}
    \begin{itemize}
        \item \textbf{前端：} Vue.js + Element Plus（国内教程丰富、组件库完善）
        \item \textbf{后端：} SpringBoot（生态成熟、易于集成开源模块）+ FastAPI（AI Agent模块轻量化）
        \item \textbf{数据库：} MySQL（团队熟悉、命令简洁）
        \item \textbf{架构：} 微服务架构（模块解耦、独立部署、故障隔离）
    \end{itemize}

    \item \textbf{核心模块与职责：}
    \begin{itemize}
        \item \textbf{课程模块：} 课程CRUD、分类管理、进度跟踪
        \item \textbf{作业模块：} 作业发布、提交、批改、成绩管理
        \item \textbf{用户模块：} 多角色（学生/教师/管理员）认证授权
        \item \textbf{AI Agent模块：} 基于FastAPI构建智能问答服务，采用检索增强生成（RAG）架构
    \end{itemize}

    \item \textbf{系统架构与部署：}
    \begin{itemize}
        \item 微服务拆分部署于4台服务器分布式系统
        \item Docker容器化部署，Nginx负载均衡与反向代理
        \item Nacos服务注册与发现，Jenkins CI/CD自动化流水线
        \item Redis缓存常用查询数据，提升系统响应速度
    \end{itemize}

    \item \textbf{开发流程：} 采用Scrum敏捷开发，定期冲刺规划、每日站会、代码审查
\end{itemize}




% 专业技能   这个真的要好好写  
\section{\makebox[\widthof{\faGraduationCap}][c]{\color{CVBlue}\faWrench}\ 专业技能}

	% \item 数据库
	% \item Linux
 %    \item shell
	% \item Java
	% \item Redis
 %    \item 常用框架
 %    \item 计算机网络 
 %    \item 虚拟化： 
 %    \item 团队协作
 %    \item 还有啥 
\begin{itemize}
    \item \textbf{数据库：} 熟悉MySQL索引原理，事务隔离级别和ACID实现机制，共享锁和排他锁机制。了解Redis底层数据结构和分布式策略。
    \item \textbf{Web框架：} 熟悉SpringBoot/MyBatis/Maven等Java开发框架的使用，了解IOC原理和AOP编程范式。
    \item \textbf{开发工具：} 熟悉Git工作流，熟悉使用JetBrain/VScode等IDE。
    \item \textbf{计算机网络：} 了解HTTP/TCP/IP/RPC等协议内容，动手实现过TCP和RPC协议客户端和服务端，了解DNS/ARP/
    DHCP/NAT等技术原理。
    \item \textbf{虚拟化：} 了解docker和docker compose的基本使用

\end{itemize}




% 其他亮点 
\section{\makebox[\widthof{\faGraduationCap}][c]{\color{CVBlue}\faTags}\ 其他} 
% increase linespacing [parsep=0.5ex]
\begin{itemize}
    % \item 对软件设计模式感兴趣，已有仓库用Java实现了23种设计模式的样例代码。
    \item 对操作系统感兴趣，完成了MIT S6.081(xv6)操作系统Labs。
\end{itemize}

 



%%%% 如果多页简历，可以手动在适当位置插入 \newpage 或者 \clearpage 开始新一页

\end{document}
